\chapter{模型}
\section{BPR模型}
% 公式一：模型打分
% 公式一：模型打分
\begin{dxtips}
   
双塔模型训练模型是为了损失函数最小，也就是所有第 $u$ 个人对第 $i$ 件事比第 $j$ 件事更喜欢的概率，因为加了负号所以他越大损失函数越小。本来应该是概率相乘但是因为太小了就取 $\ln$ 变相加了。由于概率不好算直接拿 Sigmoid 函数近似了，近似的方法就是用 $\hat{y}_{ui} - \hat{y}_{uj}$ 的差带入 Sigmoid 函数。$\hat{y}_{ui}$ 和 $\hat{y}_{uj}$ 分别是 UserTower 和 ItemTower 带入用一个 $x_u$ 再分别带入物品 $i$ 和物品 $j$ 得出的分数。

\end{dxtips}
\begin{definition}[模型打分]
$$
\hat{y}_{ui} = f_\theta(x_u)^\top g_\phi(x_i)
$$
其中 $\hat{y}_{ui}$ 为用户 $u$ 对物品 $i$ 的预测偏好分数；
$f_\theta(\cdot)$ 为 User Tower，参数为 $\theta$，输入用户特征 $x_u$，输出用户向量；
$g_\phi(\cdot)$ 为 Item Tower，参数为 $\phi$，输入物品特征 $x_i$，输出物品向量；
$\top$ 表示向量内积。
\end{definition}
\begin{note}
训练阶段整个 batch 以矩阵形式输入：
$$
X_U \in \mathbb{R}^{B \times d_u}, \quad X_V \in \mathbb{R}^{B \times d_v}
$$
其中 $B$ 为 batch size，每一行对应一个用户/物品的特征向量。

单个样本为向量形式：
$$
x_u \in \mathbb{R}^{d_u}, \quad x_i \in \mathbb{R}^{d_v}
$$
公式中 $f_\theta(x_u)$ 严格指单个样本的向量输入。
\end{note}
% 公式二：贝叶斯后验概率
\begin{note}
    embedding是嵌入的意思，说的很好transformer架构用attention公式就是把用户信息和用户的兴趣偏好嵌入了一个矩阵
\end{note}
\begin{note}
$f_\theta$ 由多层堆叠构成，上一层输出作为下一层输入：
$$
x_u \to h_1 \to h_2 \to h_3 \to \mathbf{u}
$$
每一层做一次：RELU就是保留正数部分负数部分归零
$$
h_l = \text{ReLU}(W_l h_{l-1} + b_l)
$$
最终输出即为用户 embedding $\mathbf{u} \in \mathbb{R}^d$。
\end{note}

\begin{definition}[BPR损失函数]
$$
\mathcal{L}_{\text{BPR}} = -\sum_{(u,i,j) \in \mathcal{D}} \log \sigma(\hat{y}_{ui} - \hat{y}_{uj}) + \lambda \|\Theta\|^2
$$
其中 $\mathcal{D}$ 为训练三元组集合，每条样本为 $(u, i, j)$，
即用户 $u$、正样本 $i$、负样本 $j$；
$\log \sigma(\hat{y}_{ui} - \hat{y}_{uj})$ 为对数似然，分数差越大则损失越小；
$\lambda \|\Theta\|^2$ 为 L2 正则项，$\lambda$ 控制正则化强度，防止过拟合。
\end{definition}
\begin{dxtips}
    训练是为了损失函数最小
\end{dxtips}
\begin{example}
假设模型只有一层，参数为：
$$
W = \begin{bmatrix} 1 & 2 \\ 3 & 4 \end{bmatrix}, \quad b = \begin{bmatrix} 1 \\ -1 \end{bmatrix}
$$
则 L2 正则项为：
$$
\|\Theta\|^2 = 1^2 + 2^2 + 3^2 + 4^2 + 1^2 + (-1)^2 = 1+4+9+16+1+1 = 32
$$
即将 $W$ 和 $b$ 中所有元素平方后求和，得到一个标量。
\end{example}
参数越多正则项越大。
\begin{note}
原始目标是最大化所有三元组的概率之积（连乘），但连乘结果数值极小，计算机无法存储，因此取对数将连乘变为连加：
$$
\prod_{(u,i,j)} P(i >_u j \mid \Theta) \xrightarrow{\log} \sum_{(u,i,j)} \log P(i >_u j \mid \Theta)
$$
再加负号转为最小化损失，数值稳定，结果与原始目标等价。
\end{note}
\begin{note}
每个三元组 $(u,i,j)$ 对应一个事件：用户 $u$ 更喜欢 $i$ 而不是 $j$。
训练数据中有很多这样的三元组，我们希望模型对所有这些事件同时成立的概率最大，
因此对所有三元组的概率连乘——这正是最大似然估计的思想。
\end{note}
\begin{note}
$\hat{y}_{ui} - \hat{y}_{uj}$ 越大，$\sigma$ 越接近 $1$，$\log\sigma$ 越接近 $0$，加负号后损失越小。
即正样本分数远高于负样本时，模型损失趋近于 $0$，说明排序正确。
\end{note}

\begin{definition}[贝叶斯后验概率近似方法]
$$
P(i >_u j \mid \Theta) = \sigma(\hat{y}_{ui} - \hat{y}_{uj})
$$
其中 $i >_u j$ 表示用户 $u$ 对物品 $i$ 的偏好强于物品 $j$；
$\Theta$ 为模型全部参数；
$\hat{y}_{ui}$ 与 $\hat{y}_{uj}$ 分别为用户 $u$ 对正样本 $i$（已点击）和负样本 $j$（未点击）的预测分数；
$\sigma(\cdot) = \dfrac{1}{1+e^{-x}}$ 为 Sigmoid 函数，将分数差压缩至 $(0,1)$ 区间作为概率。
\end{definition}

% 公式三：BPR损失函数
\begin{note}
贝叶斯后验概率 $P(i >_u j \mid \Theta) = \sigma(\hat{y}_{ui} - \hat{y}_{uj})$ 并非推导所得，
而是 BPR 作者的建模假设：直接用 Sigmoid 函数将正负样本的分数差映射为概率。就是为了把他映射到0到1之间
\end{note}
\begin{note}
    $\Theta$是全部样本空间是一个集合
\end{note}
\subsection{WBPR}
\begin{definition}[WBPR损失函数]
$$
\mathcal{L}_{\text{WBPR}} = -\sum_{(u,i,j) \in \mathcal{D}} w_j \cdot \log \sigma(\hat{y}_{ui} - \hat{y}_{uj}) + \lambda\|\Theta\|^2
$$
其中 $w_j$ 为负样本 $j$ 的流行度权重，越流行的 item 未被点击则权重越大，损失越大，梯度信号越强。
\end{definition}
就是在BPR的基础上给j加了权重
\chapter{ABtest}
\section{分流}
\begin{theorem}[用户级分流]
用户级分流是指以 $user\_id$ 作为实验单位，将用户稳定分配到对照组或实验组。

在推荐、增长、Feed 流、Push 等长期影响用户体验的场景中，通常优先使用用户级分流。

其核心原因是：同一个用户应当在整个实验周期内持续接受同一种策略，避免跨组污染。
\[
group_i =
\begin{cases}
control, & hash(user_i) \bmod 100 < 50, \\
treatment, & hash(user_i) \bmod 100 \geq 50.
\end{cases}
\]
\end{theorem}
\begin{note}
    哈希分流能保证稳定和随机，但有限样本下仍可能出现关键人群比例不均衡，比如实验组高活跃用户偏多。这个时候实验差异可能混入用户结构差异。所以我会在实验前后检查 SRM 和关键协变量均衡，比如新老用户、活跃度、设备、地域。如果某些关键变量很重要，可以进一步用分层随机分流，或者在分析阶段做分层分析和协变量调整。
\end{note}
\begin{dxtips}
    我们的哈希分流的缺点是不是就是可能随机分的但是分的时候某一个类型的人偏多。哈希分流的优点是可复现、可溯源、近似随机，而且能保证同一个用户稳定在同一组。对于模型对比来说，这很重要，因为我们希望控制组持续使用旧模型，实验组持续使用新模型，避免同一个用户来回切换策略造成污染。
\end{dxtips}
\begin{theorem}[设备级分流]
设备级分流是指以 $device\_id$ 作为实验单位，将设备分配到不同实验组。

这种方式适合未登录用户、新用户冷启动、App 端功能实验等场景。

其局限是：同一用户可能拥有多个设备，因此可能在不同设备上进入不同实验组。
\[
group_i = f(device\_id_i)
\]
\end{theorem}
\begin{itemize}
    \item \textbf{多设备导致跨组污染：}
    同一个真实用户可能拥有多个设备。如果按 \texttt{device\_id} 分流，不同设备可能进入不同实验组，导致同一用户同时体验对照组和实验组策略。

    \item \textbf{设备类型代表用户画像差异：}
    设备不是中性的实验单位。例如 iPhone 用户可能在收入、消费能力、城市线级、付费意愿等方面与其他用户不同。如果两组设备结构不均衡，指标差异可能混入用户画像差异。

    \item \textbf{设备性能直接影响体验指标：}
    低端机可能更卡、加载更慢、崩溃率更高，高端机体验更流畅。因此完播率、停留时长、转化率等指标可能受到设备性能影响，而不完全是策略效果。

    \item \textbf{设备 ID 不够稳定：}
    用户重装 App、清缓存、恢复出厂设置或系统隐私限制，都可能导致设备标识变化。同一设备可能被重新识别并重新分流，破坏实验稳定性。
\end{itemize}
\begin{theorem}[会话级分流]
会话级分流是指以 $session\_id$ 作为实验单位，将一次访问或一次会话分配到实验组。

这种方式适合搜索结果页、单次访问体验、临时页面策略等短周期场景。

其局限是：同一用户在不同会话中可能看到不同策略，因此不适合长期推荐策略实验。
\[
group_i = f(session\_id_i)
\]
\end{theorem}

\begin{theorem}[曝光级分流]
曝光级分流是指以 $impression\_id$ 或 $request\_id$ 作为实验单位，对每一次曝光或请求进行分组。

这种方式适合广告创意轮播、单次请求级排序实验、极短期策略测试。

其风险是：同一用户可能同时接触对照组和实验组策略，导致用户体验污染；同时曝光样本之间并不完全独立，容易造成统计显著性被高估。
\[
group_i = f(impression\_id_i)
\]
\end{theorem}

\begin{theorem}[地域级分流]
地域级分流是指以城市、省份、国家等地理单位作为实验单位。

这种方式适合区域运营、本地生活、物流履约、城市活动或政策类实验。

其局限是：地域之间本身差异较大，不一定满足随机性；同时样本单位较少，容易受到节假日、天气、城市事件等外部因素影响。
\[
group_i = f(region_i)
\]
\end{theorem}

\begin{theorem}[集群级分流]
集群级分流是指以学校、公司、社区、社交网络簇、创作者群体等 cluster 作为实验单位。

这种方式适合存在强网络效应或强干扰效应的实验，例如社交关系链、内容生态、创作者侧分发策略。

其优点是可以减少不同实验组之间的相互影响；缺点是有效样本量下降、方差变大、统计分析更复杂。
\[
group_i = f(cluster_i)
\]
\end{theorem}

\begin{theorem}[随机分流]
随机分流是指直接用随机数将实验单位分配到不同实验组。

如果实验单位为用户，则可以写作：
\[
group_i =
\begin{cases}
control, & U_i < p, \\
treatment, & U_i \geq p,
\end{cases}
\quad U_i \sim Uniform(0,1)
\]

其中 $p$ 是对照组流量比例。

随机分流的优点是简单直接；缺点是如果没有稳定键，同一个用户在不同时间可能被分到不同组。
\end{theorem}

\begin{theorem}[哈希分流]
哈希分流是指使用实验单位 ID 和实验 ID 计算哈希值，再根据哈希桶分配实验组。

工业界常用形式为：
\[
bucket_i = hash(user_i, experiment\_id) \bmod 100
\]

例如：
\[
group_i =
\begin{cases}
control, & 0 \leq bucket_i < 50, \\
treatment, & 50 \leq bucket_i < 100.
\end{cases}
\]

哈希分流的优点是稳定、可复现、易于扩展到灰度比例，并且适合多实验并行管理。
\end{theorem}

\begin{theorem}[分层随机分流]
分层随机分流是指先按照关键特征将样本划分为若干层，再在每一层内部进行随机分流。

例如按照用户活跃度 $S_i$ 分层：
\[
S_i \in \{low, medium, high\}
\]

然后在每个分层 $s$ 内部分流：
\[
P(group_i = treatment \mid S_i = s) = p
\]

这种方式可以保证新老用户、高低活跃用户、不同端类型等关键人群在实验组和对照组中更加均衡，从而降低方差，提高实验可信度。
\end{theorem}

\begin{theorem}[灰度分流]
灰度分流是指先让一小部分流量进入实验组，观察核心指标和护栏指标稳定后，再逐步扩大流量比例。

常见节奏为：
\[
1\% \rightarrow 5\% \rightarrow 10\% \rightarrow 30\% \rightarrow 50\%
\]

灰度分流适合风险较高的新功能、新模型或新推荐策略。它的目的不是单纯估计效果，而是控制上线风险。
\end{theorem}

\begin{theorem}[白名单分流]
白名单分流是指指定一批用户、员工、测试账号或特定人群进入实验组。

设白名单集合为 $W$，则：
\[
group_i =
\begin{cases}
treatment, & user_i \in W, \\
control, & user_i \notin W.
\end{cases}
\]

白名单分流适合内部测试、功能冒烟、小范围验证，但由于样本不是随机抽取，通常不能作为正式 ABTest 效果评估的依据。
\end{theorem}






\begin{definition}[用户级指标聚合]
用户级指标聚合是指在分流后，将原始行为明细按
\[
(group\_name,\ user\_id)
\]
聚合，把多条播放记录转化为每个用户一行的实验指标。

设用户 $i$ 的行为记录集合为 $\mathcal{D}_i$，则用户级指标为：
\[
M_i = f(\mathcal{D}_i)
\]

例如用户完播率：
\[
completion\_rate_i
=
\frac{
\sum_{j \in \mathcal{D}_i}
\mathbf{1}(play\_duration_{ij} \geq video\_duration_{ij})
}{
\sum_{j \in \mathcal{D}_i}
\mathbf{1}(play\_duration_{ij} > 0)
}
\]

其目的，是让每个用户作为一个实验样本，再比较实验组与对照组的平均表现。
\end{definition}
\begin{example}[两层指标聚合]
本项目的 ABTest 指标聚合分为两层。

第一层，先按实验组和用户聚合：
\[
GROUP\ BY\ (group\_name,\ user\_id)
\]

例如，对每个用户计算完播率：
\[
completion\_rate_i
=
\frac{complete\_cnt_i}{play\_cnt_i}
\]

第二层，再按实验组计算组内平均：
\[
\bar{M}_{g}
=
\frac{1}{n_g}
\sum_{i:g_i=g} M_i,
\quad
g \in \{control,\ treatment\}
\]

最终比较实验组和对照组的均值差异：
\[
\bar{M}_{treatment} - \bar{M}_{control}
\]

也就是：先给每个用户算成绩，再比较两组用户的平均成绩。
\end{example}
\begin{itemize}
    \item \textbf{指标层级要和实验单位一致：}
    本项目按 \texttt{user\_id} 分流，因此后续也要聚合到用户级：
    \[
    GROUP\ BY\ (group\_name,\ user\_id)
    \]
    避免直接用曝光级数据导致高活跃用户被重复放大。

    \item \textbf{完播率先按用户计算，再按组求平均：}
    先计算每个用户的完播率：
    \[
    completion\_rate_i
    =
    \frac{complete\_cnt_i}{play\_cnt_i}
    \]
    再比较两组用户级均值：
    \[
    \bar{M}_{treatment} - \bar{M}_{control}
    \]
    这样每个用户的权重更接近一致。

    \item \textbf{处理分母为 0 的情况：}
    当用户没有播放行为时，完播率分母可能为 $0$。项目中使用：
    \[
    complete\_cnt / NULLIF(play\_cnt, 0)
    \]
    并通过 \texttt{WHERE play\_cnt > 0} 保留有效用户。

    \item \textbf{处理异常观看比：}
    对 \texttt{watch\_ratio} 做截断：
    \[
    watch\_ratio \in [0,1]
    \]
    避免异常值拉偏平均观看比。

    \item \textbf{保留分层变量：}
    聚合时保留 \texttt{user\_active\_degree}，用于后续按用户活跃度做分层分析，判断实验效果是否在不同用户群体上一致。
\end{itemize}
 \begin{example}[曝光级重复放大问题]
如果直接使用曝光级数据进行比较，高活跃用户会因为产生更多曝光记录而被重复计入样本。

例如，用户 $i$ 产生 $100$ 条曝光，用户 $k$ 只产生 $1$ 条曝光，则在曝光级分析中，用户 $i$ 的权重约为用户 $k$ 的 $100$ 倍。

但在用户级分流的 ABTest 中，实验单位是用户。因此，更合理的做法是先将曝光记录聚合到用户级，使每个用户先形成一个指标：
\[
M_i = f(\mathcal{D}_i)
\]

再比较实验组和对照组的用户级平均表现：
\[
\bar{M}_{treatment} - \bar{M}_{control}
\]

这样可以避免高活跃用户因曝光次数更多而过度影响实验结论。
\end{example}
\begin{dxtips}
    这里假设的就是H0 μ1-μ2=0
    也就是正态分布以原点为对称轴，用实际的 $
    \bar X_1-\bar X_2
    $算出后看看落在正态分布的那个点上，这个其实能第一看看在不在拒绝域，第二能算p值看看是不是显著。用体验是总体方差不知道所以用样本方差s，(n-1)s/σ²符合卡方分布，自由度是n-1，然后变成标准卡方分布就又把n-1干掉了。上面是标准正态分布。
\end{dxtips}
\begin{enumerate}
    \item \textbf{只看 $p$ 值，不看业务大小}

    样本量很大时，极小的差异也可能显著。例如完播率只提升 $0.01\%$，也可能得到
    \[
    p < 0.05
    \]
    但这不一定有业务价值。因此统计检验后还要看 lift、置信区间和业务成本。

    \item \textbf{曝光级样本不独立}

    如果直接使用曝光级数据做检验，高活跃用户会贡献多条曝光记录，导致样本之间不独立，显著性可能被高估。

    因此，本项目更重视用户级指标上的 t-test。

    \item \textbf{多指标重复检验}

    ABTest 往往同时观察多个指标，例如完播率、观看时长、点赞率、留存率、负反馈率等。

    指标越多，偶然显著的概率越高。因此需要提前确定核心指标，其他指标作为辅助；必要时做多重检验校正。

    \item \textbf{方差不齐或分布偏态}

    用户行为指标可能存在长尾，例如少数用户观看时长特别高，导致均值和方差被拉动。

    因此，Welch t-test 比普通 t-test 更稳健；如果分布极度偏态，也可以考虑 Bootstrap 或非参数检验。

    \item \textbf{分流异常或实验污染}

    如果分流本身有问题，后续检验就不可靠。例如实验组高活跃用户明显更多、两组人数比例异常、同一用户跨组、埋点口径不一致等。

    因此，在正式解释检验结果前，需要先检查 SRM 和关键特征均衡。
\end{enumerate}
\begin{definition}[Welch t-test]
Welch t-test 是一种用于比较两个独立样本均值是否相等的双样本 $t$ 检验。与普通双样本 $t$ 检验不同，Welch t-test 不要求两组总体方差相等，即允许 $\sigma_1^2 \neq \sigma_2^2$。

设两组样本均值、样本方差和样本量分别为 $\bar X_1,s_1^2,n_1$ 与 $\bar X_2,s_2^2,n_2$，其检验统计量为
\[
t=\frac{\bar X_1-\bar X_2}{\sqrt{\frac{s_1^2}{n_1}+\frac{s_2^2}{n_2}}}.
\]

因此，Welch t-test 适用于两组样本方差可能不相等的均值比较问题。
\end{definition}
\begin{note}[分层分析与辛普森悖论]
分层分析的目的，是检查整体实验结论是否在不同用户群体中仍然成立。

若只看整体均值，可能出现辛普森悖论：整体上实验组优于对照组，但在某些分层中实验组并不占优，甚至更差。

因此，在 ABTest 中需要按关键变量（如用户活跃度、新老用户、设备类型等）做分层分析，判断策略效果是否稳定，以及是否伤害了某些用户群体。
\end{note}
\begin{example}[分层分析中的异质性效果]
整体结果显示：
\[
\bar{M}_{\text{实}} > \bar{M}_{\text{对}}
\]
似乎说明实验组策略优于对照组。

但按用户活跃度分层后，可能得到：
\[
\begin{array}{c|c}
\text{用户分层} & \text{实验效果} \\
\hline
\text{高活跃用户} & \text{实验组更好} \\
\text{中活跃用户} & \text{实验组与对照组差不多} \\
\text{低活跃用户} & \text{实验组更差}
\end{array}
\]

这说明整体提升可能主要来自高活跃用户，而低活跃用户反而被新策略伤害。因此，ABTest 不能只看整体均值，还需要结合分层结果判断策略是否稳定。
\end{example}
\begin{definition}[ABTest 结果解读]
ABTest 结果解读是指在完成统计检验和分层分析后，将统计结果转化为业务结论。

结果解读主要关注四个方面：
\begin{enumerate}
    \item \textbf{方向：}实验组指标是高于还是低于对照组，即判断
    \[
    \bar{M}_{\text{实}}-\bar{M}_{\text{对}}
    \]
    的正负。

    \item \textbf{显著性：}差异是否达到统计显著，例如是否满足
    \[
    p < 0.05.
    \]

    \item \textbf{业务大小：}即使统计显著，也要判断 lift 是否足够大，是否具有实际业务价值。

    \item \textbf{分层稳定性：}检查不同用户分层中的效果是否一致，避免整体提升但某些关键人群受损。
\end{enumerate}

因此，结果解读的目的不是单纯判断 $p$ 值，而是回答：实验是否真的带来了值得推进的业务提升。
\end{definition}